%%%%%%%%%%%%%%%%%%%%%%%%%%%%%%%%
%%%
%%% Research Summary 
%%%
%%% Author - Steve Hurder
%%%
%%% Date Started: October 12, 2009
%%% Date Completed: November 15 , 2009
%%%%%%%%%%%%%%%%%%%%%%%%%%%%%%%%

\documentclass[12pt]{amsart}
\usepackage{graphicx,wrapfig}
\usepackage{amssymb}
\usepackage{epstopdf}
\usepackage{nicefrac}
\DeclareGraphicsRule{.tif}{png}{.png}{`convert #1 `dirname #1`/`basename #1 .tif`.png}

\parskip = 6pt
\parindent = 0.0in

\hoffset=-1in
\voffset=-.7in
\setlength{\textwidth}{7in}
\setlength{\textheight}{9in}

%%%%%%%%%%%%%%%%%%%%%%%%%%%%%%%%%%%%%%%
%%%%%%%%%%%% Preamble %%%%%%%%%%%%%%%%%
%%%%%%%%%%%%%%%%%%%%%%%%%%%%%%%%%%%%%%%
%% Flexible formatting options
%\usepackage[margin=1.25in]{geometry} %% Change margins

\usepackage{xcolor} % Required for inserting images
\usepackage{amsthm,,amssymb}
\usepackage{enumerate,enumitem,nicefrac,listings,longtable,caption,makecell}
% Make references linky
\usepackage[pdfencoding=auto,psdextra,colorlinks,linkcolor=blue,citecolor=blue]{hyperref}
% Insert line numbers
\usepackage[pagewise]{lineno}\nolinenumbers %no line numbers

%% Theorem environments
\newtheorem{theorem}{Theorem}[section]
\theoremstyle{plain}
\newtheorem{lemma}[theorem]{Lemma}
\newtheorem{conjecture}[theorem]{Conjecture}
\newtheorem{proposition}[theorem]{Proposition}
\newtheorem{definition}[theorem]{Definition}
\newtheorem{corollary}[theorem]{Corollary}
\theoremstyle{definition}
\newtheorem{remark}[theorem]{Remark}
\newtheorem{example}[theorem]{Example}

%% Math macros
\newcommand{\q}[1]{\overline{#1}}
\newcommand{\mc}[1]{\mathcal{#1}}
\newcommand{\mbf}[1]{\mathbf{#1}}
\newcommand{\mr}[1]{\mathrm{#1}}
\newcommand{\gen}[1]{\langle {#1} \rangle}
\newcommand{\nrm}{\trianglelefteq}
\newcommand{\N}{\mathbb{N}} %Bold N, naturals
\newcommand{\C}{\mathbb{C}}
\newcommand{\Zen}[1]{\mathbf{Z}(#1)} %center
\newcommand{\Irr}[1]{\mathrm{Irr}(#1)} 
\newcommand{\Cen}[2]{\mathbf{C}_{#1}(#2)} % centralizer in #1 of #2
\newcommand{\Norm}[2]{\mathbf{N}_{#1}(#2)} % normalizer in #1 of #2
\renewcommand{\char}{\, \textrm{char} \,} % characteristic
\newcommand{\wh}{\widehat}

%% Groups of Lie Type
\renewcommand{\sl}[2]{\mathrm{SL}_{#1}(#2)} % SL_{#1}(#2)
\newcommand{\pgl}[2]{\mathrm{PGL}_{#1}(#2)} % PGL_{#1}(#2)
\newcommand{\psl}[2]{\mathrm{PSL}_{#1}(#2)} % PSL_{#1}(#2)
\newcommand{\su}[2]{\mathrm{SU}_{#1}(#2)}   % SU_{#1}(#2)
\newcommand{\psu}[2]{\mathrm{PSU}_{#1}(#2)} % PSU_{#1}(#2)
\renewcommand{\sp}[2]{\mathrm{Sp}_{#1}(#2)} % Sp_{#1}(#2)
\newcommand{\psp}[2]{\mathrm{PSp}_{#1}(#2)} % PSp_{#1}(#2)



\begin{document}

 
\title{Finite Groups I: List of Errata}

%\thanks{2000 {\it Mathematics Subject Classification}. Primary 20E34, 20C15; Secondary 81T45, 20E35}
 
%\author{Chris Schroeder}
%\address{Department of Mathematics and Statistics, Binghamton University,
%    Binghamton, NY 13902-6000, USA}
%\email{cschroe2@binghamton.edu}
\thanks{Last updated: \today}

%\date{\today}
 

%\keywords{}

\maketitle

%\vspace{-30pt}

\begin{enumerate}[label=\arabic*.]
    \item pg.~223, Main Theorem II.8.27, item (7): 
    
    The 1979 errata, which are included in the English translation, list the conditions on $t$ as $t \mid (p^m-1)/d$ and $t \mid (p^f-1)$. But this excludes subgroups that do exist. For example, $\mr{PSL}_2(7^4)$ contains subgroups of order $6\cdot 7^3$. (Thank you to I.~Heckenberger for this example.) As the statement is not an ``if and only if'', the 1967 German original is technically correct. In the original, the conditions on $t$ are given as $t \mid (p^m-1)$ and $t \mid (p^f-1)$. Note that this allows subgroups that do not exist. For example, when $m=f$, we know that we can only have $p^f:(p^f-1)/2$.  Or similarly, $\mr{PSL}_2(27)$ would have a subgroup $3:2\cong S_3$, but there is no such subgroup. (Thank you to C.~Parker for this example.) A more precise condition would be $t \mid (p^m-1)$ and $t \mid (p^f-1)/d$, although this still allows too many semidirect products that cannot be realized as subgroups.  Thank you to I.~Heckenberger, G.~Malle, C.~Parker, B.~Sambale, H.~Tong-Viet and R.~Waldecker for addressing this error.
\end{enumerate}

\end{document}